\documentclass[11pt,a4paper]{article}
\usepackage[utf8]{inputenc}
\usepackage[T1]{fontenc}
\usepackage[ngerman]{babel}
\usepackage{helvet}
\renewcommand{\familydefault}{\sfdefault}
\usepackage[left=2cm, right=2cm, top=2cm, bottom=2cm]{geometry}
\usepackage{amsmath,amssymb}
\usepackage{array}
\usepackage{booktabs}
\usepackage{tikz}
\usepackage{pgfplots}
\pgfplotsset{compat=1.18}
\usepackage{fancyhdr}

\pagestyle{fancy}
\fancyhf{}
\fancyhead[L]{\small Physik Klasse 10E}
\fancyhead[R]{\small Probeklausur Kinematik -- Variante B}
\fancyfoot[C]{\small Seite \thepage}
\renewcommand{\headrulewidth}{0.4pt}

\setlength{\parindent}{0pt}
\setlength{\fboxsep}{8pt}

% Aufgaben-Befehl im Layout B Stil
\newcounter{aufgabe}
\newcommand{\aufgabe}[2]{%
\stepcounter{aufgabe}%
\vspace{0.5cm}%
\noindent\fbox{\parbox{\dimexpr\textwidth-2\fboxsep-2\fboxrule}{%
\textbf{Aufgabe \theaufgabe: #1}\hfill\textbf{#2 BE}%
}}%
\vspace{0.3cm}%
}

\begin{document}

% === KOPFBEREICH ===
\noindent\fbox{\parbox{\dimexpr\textwidth-2\fboxsep-2\fboxrule}{
\centering
{\Large\bfseries Probeklausur: Kinematik -- Variante B}\\[0.2cm]
{\small Klasse 10E \quad|\quad 90 Minuten \quad|\quad 60 BE}
}}

\vspace{0.4cm}
\begin{tabular}{@{}p{6cm}p{8cm}@{}}
Name: & \hrulefill \\[0.3cm]
Datum: & \hrulefill \\
\end{tabular}

\vspace{0.3cm}

\begin{center}
\small
\begin{tabular}{|c|c|c|c|c|c|c|c|c|c|}
\hline
\textbf{Aufgabe} & 1 & 2 & 3 & 4 & 5 & 6 & 7 & Zusatz & \textbf{Summe} \\
\hline
\textbf{max. BE} & 6 & 8 & 8 & 8 & 8 & 10 & 10 & (2) & 60 \\
\hline
\textbf{erreicht} & & & & & & & & & \\
\hline
\end{tabular}
\end{center}

\vspace{0.3cm}

% === AUFGABEN ===

\aufgabe{Grundlagen}{6}

\begin{enumerate}
\item[a)] Rechne um: 90 km/h = \underline{\hspace{2cm}} m/s \hfill (1 BE)

\item[b)] Ein Stein fällt von einer Brücke (Höhe 20 m). Berechne die Fallzeit bis zum Aufprall auf dem Wasser. \hfill (2 BE)

\vspace{2.5cm}

\item[c)] Ergänze die Tabelle: \hfill (3 BE)

\begin{center}
\begin{tabular}{|l|c|c|}
\hline
\textbf{Diagramm} & \textbf{Steigung bedeutet} & \textbf{Fläche bedeutet} \\
\hline
$s(t)$-Diagramm & & --- \\
\hline
$v(t)$-Diagramm & & \\
\hline
\end{tabular}
\end{center}
\end{enumerate}

\aufgabe{Notbremsung auf der Autobahn}{8}

Ein LKW fährt mit 90 km/h auf der Autobahn. Plötzlich entsteht ein Stauende. Der Fahrer bremst mit $a = -5\,\mathrm{m/s^2}$.

\begin{enumerate}
\item[a)] Rechne die Geschwindigkeit in m/s um. \hfill (1 BE)

\vspace{1.2cm}

\item[b)] Berechne die Bremszeit bis zum Stillstand. \hfill (2 BE)

\vspace{2cm}

\item[c)] Berechne den Bremsweg. \hfill (2 BE)

\vspace{2.5cm}

\item[d)] Das Stauende ist 50 m entfernt. Kann der LKW rechtzeitig anhalten? Begründe und berechne den Abstand zum Stauende. \hfill (3 BE)

\vspace{2.5cm}
\end{enumerate}

\newpage

\aufgabe{Fußball-Abstoß}{8}

Ein Torwart schießt den Ball senkrecht nach oben mit $v_0 = 20\,\mathrm{m/s}$.

\begin{enumerate}
\item[a)] Berechne die Steigzeit (Zeit bis zum höchsten Punkt). \hfill (2 BE)

\vspace{2.5cm}

\item[b)] Berechne die maximale Steighöhe über der Abschussposition. \hfill (3 BE)

\vspace{3cm}

\item[c)] Mit welcher Geschwindigkeit kommt der Ball wieder auf Abschusshöhe an? Begründe ohne Rechnung. \hfill (3 BE)

\vspace{2.5cm}
\end{enumerate}

\aufgabe{Fahrradfahrt}{8}

Das $v(t)$-Diagramm zeigt die Fahrt eines Fahrrads:

\begin{center}
\begin{tikzpicture}
\begin{axis}[
    width=11cm, height=5cm,
    xlabel={$t$ in s},
    ylabel={$v$ in m/s},
    xmin=0, xmax=15,
    ymin=0, ymax=12,
    grid=major,
    axis lines=left,
    xtick={0,3,6,9,12,15},
    ytick={0,2,4,6,8,10,12}
]
% Schattierte Flächen zur Visualisierung
\addplot[fill=black!10, draw=none] coordinates {(0,0) (6,0) (6,6) (0,0)} \closedcycle;
\addplot[fill=black!10, draw=none] coordinates {(6,0) (12,0) (12,6) (6,6)} \closedcycle;
\addplot[fill=black!10, draw=none] coordinates {(12,0) (15,0) (12,6)} \closedcycle;
% Graph
\addplot[thick, black] coordinates {(0,0) (6,6) (12,6) (15,0)};
% Phasen-Beschriftungen
\node at (axis cs:3,2) {\small A};
\node at (axis cs:9,3) {\small B};
\node at (axis cs:13.5,2) {\small C};
\end{axis}
\end{tikzpicture}
\end{center}

\begin{enumerate}
\item[a)] Beschreibe die drei Phasen A, B und C in Worten. \hfill (3 BE)

\vspace{2.5cm}

\item[b)] Berechne die Beschleunigung in Phase A. \hfill (2 BE)

\vspace{2cm}

\item[c)] Berechne den Gesamtweg der 15-sekündigen Fahrt. \hfill (3 BE)

\vspace{3cm}
\end{enumerate}

\newpage

\aufgabe{Fallschirm und Luftwiderstand}{8}

\begin{enumerate}
\item[a)] Ein Fallschirmspringer springt aus einem Flugzeug. Zunächst beschleunigt er, dann fällt er mit konstanter Geschwindigkeit. Erkläre physikalisch, warum er nicht immer schneller wird. Nenne zwei relevante physikalische Größen. \hfill (3 BE)

\vspace{3.5cm}

\item[b)] Ein Tischtennisball und eine Stahlkugel gleicher Größe werden aus 10 m Höhe fallen gelassen. Erkläre, warum die Stahlkugel zuerst ankommt. \hfill (3 BE)

\vspace{3.5cm}

\item[c)] 1971 ließ Astronaut David Scott auf dem Mond einen Hammer und eine Feder gleichzeitig fallen. Beide trafen gleichzeitig auf.

Beurteile die Aussage: "`Im Vakuum fallen alle Körper gleich schnell."' \hfill (2 BE)

\vspace{3cm}
\end{enumerate}

\aufgabe{Überholvorgang auf der Landstraße}{10}

Ein PKW fährt mit konstant 72 km/h hinter einem LKW (konstant 54 km/h). Der Abstand beträgt anfangs 20 m.

\begin{enumerate}
\item[a)] Rechne beide Geschwindigkeiten in m/s um. \hfill (2 BE)

\vspace{1.5cm}

\item[b)] Mit welcher Relativgeschwindigkeit nähert sich der PKW dem LKW? \hfill (2 BE)

\vspace{2cm}

\item[c)] Der PKW beginnt zu überholen. Nach welcher Zeit ist er auf gleicher Höhe mit dem LKW? \hfill (2 BE)

\vspace{2cm}

\item[d)] Welche Strecke legt der PKW in dieser Zeit zurück? \hfill (2 BE)

\vspace{2cm}

\item[e)] Erkläre, warum Überholvorgänge bei geringem Geschwindigkeitsunterschied gefährlich sind. \hfill (2 BE)

\vspace{2cm}
\end{enumerate}

\newpage

\aufgabe{Sprung vom 10-Meter-Turm}{10}

Ein Turmspringer springt vom 10-Meter-Brett senkrecht nach unten.

\begin{enumerate}
\item[a)] Berechne die Fallzeit bis zum Eintauchen ins Wasser. \hfill (2 BE)

\vspace{2.5cm}

\item[b)] Berechne die Geschwindigkeit beim Eintauchen. \hfill (2 BE)

\vspace{2.5cm}

\item[c)] Der Springer springt nun mit einer Anfangsgeschwindigkeit von $v_0 = 2\,\mathrm{m/s}$ nach unten ab. Berechne die neue Eintauchgeschwindigkeit. \hfill (3 BE)

\vspace{3cm}

\item[d)] Vergleiche die beiden Eintauchgeschwindigkeiten. Erkläre, warum der Unterschied relativ gering ist. \hfill (3 BE)

\vspace{3cm}
\end{enumerate}

\vspace{0.5cm}
\hrule
\vspace{0.3cm}

\noindent\fbox{\parbox{\dimexpr\textwidth-2\fboxsep-2\fboxrule}{%
\textbf{Zusatzaufgabe} (für schnelle Schüler, max. 2 Zusatz-BE)
}}

\vspace{0.3cm}

Ein Auto startet aus dem Stand mit $a = 2\,\mathrm{m/s^2}$. Gleichzeitig startet 25 m dahinter ein Radfahrer mit konstanter Geschwindigkeit $v = 10\,\mathrm{m/s}$.

Nach welcher Zeit holt der Radfahrer das Auto ein? \hfill (2 BE)

\vspace{3cm}

\hrule
\vspace{0.3cm}
\begin{center}
\textit{Erlaubte Hilfsmittel: Tafelwerk}\\[0.2cm]
\textit{Hinweis: $g \approx 10\,\mathrm{m/s^2}$}
\end{center}

\end{document}
